%============================================================
%  TENSOR FORMATS IN BANACH SPACES:
%  Minimal Subspaces, Geometry, and Variational Principles
%
%  Antonio Falcó
%  Universidad Cardenal Herrera-CEU, CEU Universities
%
%  Skeleton for Ergebnisse der Mathematik und ihrer
%  Grenzgebiete (3. Folge), Springer
%
%  Estimated length: ~320 pages
%============================================================

\documentclass[envcountsame,envcountchap,sectrefs]{svmono}

% ---- Standard packages ----
\usepackage[utf8]{inputenc}
\usepackage{amsmath, amssymb, amsthm, mathrsfs}
\usepackage{mathtools}          % \coloneqq, etc.
\usepackage{graphicx}
\usepackage{tikz}
\usepackage[all,cmtip]{xy}
\usepackage{synttree}           % dimension partition trees
\usepackage{hyperref}
\usepackage{makeidx}
\makeindex

%============================================================
%  THEOREM-LIKE ENVIRONMENTS
%============================================================
\newtheorem{theorem}{Theorem}[section]
\newtheorem{lemma}[theorem]{Lemma}
\newtheorem{proposition}[theorem]{Proposition}
\newtheorem{corollary}[theorem]{Corollary}

\theoremstyle{definition}
\newtheorem{definition}[theorem]{Definition}
\newtheorem{example}[theorem]{Example}
\newtheorem{remark}[theorem]{Remark}
\newtheorem{notation}[theorem]{Notation}
\newtheorem{assumption}[theorem]{Assumption}

%============================================================
%  UNIFIED NOTATION  --  ALL MACROS DEFINED HERE
%
%  PHILOSOPHY:
%  (a) Scalars/indices: lowercase Roman or Greek
%  (b) Vectors in component spaces: lowercase bold  \bv, \bu ...
%  (c) Tensors: uppercase bold  \bV, \bU ...
%  (d) Tensor spaces (algebraic):  \VD, \Valpha, ...
%  (e) Tensor Banach spaces (topological completion): with norm subscript
%  (f) Trees, partitions, index sets: calligraphic or Roman
%============================================================

% ---- Fields ----
\newcommand{\RR}{\mathbb{R}}          % real numbers
\newcommand{\CC}{\mathbb{C}}          % complex numbers
\newcommand{\KK}{\mathbb{K}}          % generic field (R or C)
\newcommand{\NN}{\mathbb{N}}          % natural numbers (positive)
\newcommand{\NNz}{\mathbb{N}_0}       % non-negative integers
\newcommand{\ZZ}{\mathbb{Z}}          % integers

% ---- Index sets ----
% D = {1,...,d}  is the dimension index set
\newcommand{\DD}{D}                    % dimension index set
\newcommand{\dD}{d}                    % cardinality of D

% ---- Vector spaces (component spaces) ----
% V_j for j in D  (scalar spaces)
% Written as V_j or V_alpha for alpha a subset of D
\newcommand{\Vj}{V_j}                  % j-th component space
\newcommand{\Valpha}{V_\alpha}         % space associated to alpha
\newcommand{\Vbeta}{V_\beta}           % space associated to beta
\newcommand{\Vnu}{V_\nu}              % space associated to nu (multi-index)

% ---- Algebraic tensor spaces ----
% Convention: _a\bigotimes denotes algebraic tensor product
% \VD  = algebraic tensor product of all V_j, j in D
\newcommand{\VD}{\mathbf{V}_\DD}
\newcommand{\VDalg}{{}_a\bigotimes_{\alpha \in \DD} V_\alpha}
\newcommand{\Vbigalpha}{\mathbf{V}_\alpha}  % V_alpha = alg. tensor over alpha
\newcommand{\Vbigbeta}{\mathbf{V}_\beta}
\newcommand{\VPk}{\mathbf{V}_{\mathcal{P}_k}} % space at level k partition

% Algebraic tensor product symbol
\newcommand{\atens}{\otimes_a}          %  v \atens w  (algebraic)
\newcommand{\Atens}{\bigotimes_a}       %  \Atens_{j} V_j

% ---- Topological tensor spaces (completions) ----
% Convention: completion of V_D under norm ||\cdot||_D
%   written as  \VDnorm  or  V_{D,\|\cdot\|_D}
\newcommand{\VDnorm}{\mathbf{V}_{D,\|\cdot\|_D}}
\newcommand{\Valphanorm}{\mathbf{V}_{\alpha,\|\cdot\|_\alpha}}
% General: _{\|\cdot\|}\bigotimes  for topological tensor product
\newcommand{\tnorm}[1]{{}_{\|\cdot\|_{#1}}\bigotimes}

% ---- Tensors (elements) ----
% Bold lowercase: component vectors
\newcommand{\bv}{\mathbf{v}}    % generic tensor
\newcommand{\bu}{\mathbf{u}}    % generic tensor
\newcommand{\bw}{\mathbf{w}}    % generic tensor
\newcommand{\bx}{\mathbf{x}}    % generic tensor
\newcommand{\bF}{\mathbf{F}}    % right-hand side in ODE
\newcommand{\bU}{\mathbf{U}}    % block tensor
% Elementary tensor
\newcommand{\etens}[1]{\bigotimes_{#1}}  % elementary tensor symbol

% ---- Norms ----
\newcommand{\norm}[2]{\left\|#1\right\|_{#2}}   % \norm{v}{D}
\newcommand{\normD}[1]{\left\|#1\right\|_{D}}   % norm in V_D
\newcommand{\norma}[1]{\left\|#1\right\|_{\alpha}}
\newcommand{\normj}[1]{\left\|#1\right\|_{j}}
\newcommand{\dualnorm}[1]{\left\|#1\right\|_{#1^*}}
\newcommand{\opnorm}[3]{\left\|#1\right\|_{#2 \leftarrow #3}} % operator norm

% ---- Dual spaces ----
\newcommand{\dual}[1]{#1^*}       % topological dual
\newcommand{\algdual}[1]{#1'}     % algebraic dual  L(V,K)

% ---- Linear maps and operators ----
\newcommand{\cL}{\mathcal{L}}     % space of bounded linear maps
\newcommand{\GL}[1]{\mathrm{GL}(#1)}  % bounded automorphisms
\newcommand{\Id}[1]{\mathrm{id}_{#1}} % identity on #1
\newcommand{\Proj}[2]{P_{#1 \oplus #2}} % projection onto #1 along #2
\newcommand{\Isom}[2]{\cL(#1,#2)} % bounded linear maps

% ---- Subspaces, Grassmannian ----
\newcommand{\Grass}[2]{\mathbb{G}_{#1}(#2)}  % Grassmannian of r-dim subspaces
\newcommand{\GrassB}[1]{\mathbb{G}(#1)}       % Grassmannian (all split subspaces)
\newcommand{\AD}[1]{\mathcal{AD}(#1)}         % set of admissible ranks

% ---- Minimal subspaces ----
% U^min_alpha(v)  =  minimal subspace of v associated to alpha
\newcommand{\Umin}[2]{U^{\min}_{#1}(#2)}         % \Umin{\alpha}{\bv}
\newcommand{\Uminj}[1]{U^{\min}_{j}(#1)}          % \Uminj{\bv}
\newcommand{\Wmin}[2]{W^{\min}_{#1}(#2)}          % complement in decomp.
\newcommand{\alphrank}[2]{r_{#1}(#2)}             % alpha-rank of v = dim Umin
% Short version when tensor is clear from context:
\newcommand{\Uminalpha}{U^{\min}_\alpha(\bv)}
\newcommand{\alphabrank}{r_\alpha}  % alpha-rank (fixed)

% ---- Tucker format ----
% M_r(V_D)  = set of tensors in Tucker format with fixed Tucker rank r
\newcommand{\Tucker}[2]{\mathfrak{M}_{#1}(#2)}   % \Tucker{\fr}{V_D}
\newcommand{\TuckerV}{\mathfrak{M}_{\fr}(\VD)}    % with canonical names
% Tangent-like subspace at v in Tucker manifold
\newcommand{\ZD}[1]{Z^{(D)}(#1)}                  % \ZD{\bv}

% ---- Tree-based format ----
% T_D  = dimension partition tree over D
\newcommand{\TD}{T_D}                  % dimension partition tree
\newcommand{\TDtuck}{T_D^{\mathrm{Tucker}}}  % Tucker tree (depth 1)
\newcommand{\TDtt}{T_D^{\mathrm{TT}}}       % Tensor Train tree
\newcommand{\TDht}{T_D^{\mathrm{HT}}}       % Hierarchical Tucker tree

% Sons (children) of a vertex alpha in T_D
\newcommand{\Sons}[1]{S(#1)}           % \Sons{\alpha}
% Leaves of T_D
\newcommand{\Leaves}{\mathcal{L}(\TD)}
% Levels and depth
\newcommand{\lev}[1]{\mathrm{level}(#1)}
\newcommand{\dep}[1]{\mathrm{depth}(#1)}
% Partition at level k induced by T_D
\newcommand{\Pk}{\mathcal{P}_k(\TD)}   % k-th level partition
\newcommand{\PkT}[2]{\mathcal{P}_{#1}(#2)}  % general: \PkT{k}{T_D}

% FT_r(V_D, T_D)  =  set of tree-based tensors with fixed rank r
\newcommand{\FT}[3]{\mathcal{FT}_{#1}(#2,#3)}  % \FT{\fr}{\VD}{\TD}
\newcommand{\FTv}{\mathcal{FT}_{\fr}(\VD,\TD)}  % canonical
\newcommand{\FTleq}[3]{\mathcal{FT}_{\leq #1}(#2,#3)}  % bounded rank

% Tree-based rank tuple
\newcommand{\fr}{\mathfrak{r}}         % tree-based rank tuple
\newcommand{\fs}{\mathfrak{s}}         % another rank tuple
% Admissible tuples for (V_D, T_D)
\newcommand{\ADTD}{\mathcal{AD}(\VD,\TD)}

% ---- Rank-related objects ----
% C^(alpha) = transfer tensor at vertex alpha
\newcommand{\Calpha}{C^{(\alpha)}}
\newcommand{\CD}{C^{(D)}}
\newcommand{\Cmu}{C^{(\mu)}}
% Matricisation map
\newcommand{\Mmat}[1]{\mathcal{M}_{#1}}    % \Mmat{\beta}(C)

% ---- Banach manifold objects ----
\newcommand{\atlas}{\mathcal{A}}       % atlas
\newcommand{\chart}{\mathcal{U}}       % chart domain
% Local chart map at v
\newcommand{\chartmap}[1]{\xi_{#1}}    % \chartmap{\bv}
\newcommand{\Psi}{\Psi_{\bv}}          % chart map (Grassmann)
% Tangent space
\newcommand{\Tv}[2]{T_{#1}#2}          % tangent space T_v M
% Model Banach space for Tucker manifold at v
\newcommand{\Emodel}[1]{\mathbf{E}_{#1}}  % model space

% ---- Dirac-Frenkel objects ----
% Projection onto tangent space
\newcommand{\Pmetric}[1]{\mathcal{P}_{#1}}   % metric projection at v
% ODE: \dot{u} = F(t,u)
\newcommand{\Fvf}{\mathbf{F}}          % vector field in ODE

% ---- Crossnorm / norm conditions ----
\newcommand{\injnorm}{\|\cdot\|_\varepsilon}   % injective norm
\newcommand{\projnorm}{\|\cdot\|_\pi}          % projective norm
% Condition (inj): norm >= injective norm
\newcommand{\condINJ}{\textup{(Inj)}}
% Condition (cross): norm is a crossnorm
\newcommand{\condCROSS}{\textup{(Cross)}}

% ---- Miscellaneous ----
\newcommand{\spann}{\mathrm{span}}      % linear span
\newcommand{\rank}{\mathrm{rank}}       % matrix rank
\newcommand{\depth}{\mathrm{depth}}     % depth of tree
\newcommand{\level}{\mathrm{level}}     % level in tree
\newcommand{\diam}{\mathrm{diam}}       % diameter
\newcommand{\dist}{\mathrm{dist}}       % distance
\newcommand{\cl}[1]{\overline{#1}}      % closure
\newcommand{\interior}[1]{\mathring{#1}} % interior
\newcommand{\card}[1]{\#(#1)}           % cardinality
% Complement of alpha in D
\newcommand{\ac}{\alpha^c}              % D \ alpha
% Power set of D
\newcommand{\pset}[1]{2^{#1}}           % 2^D

%============================================================
%  CROSS-REFERENCE SHORTHANDS
%============================================================
\newcommand{\refthm}[1]{Theorem~\ref{#1}}
\newcommand{\reflem}[1]{Lemma~\ref{#1}}
\newcommand{\refprop}[1]{Proposition~\ref{#1}}
\newcommand{\refdef}[1]{Definition~\ref{#1}}
\newcommand{\refsec}[1]{Section~\ref{#1}}
\newcommand{\refchap}[1]{Chapter~\ref{#1}}
\newcommand{\refeq}[1]{(\ref{#1})}

%============================================================
%  DOCUMENT BEGINS
%============================================================
\begin{document}

\frontmatter

\title{Tensor Formats in Banach Spaces:\\
Minimal Subspaces, Geometry,\\
and Variational Principles}

\author{Antonio Falcó}

\maketitle

%------------------------------------------------------------
\chapter*{Preface}
\label{chap:preface}
%------------------------------------------------------------
% [~8 pp.]
%
% CONTENT OUTLINE:
%
% 1. Motivation and context
%    - The curse of dimensionality in modern mathematics and applications
%    - Low-rank tensor formats as a strategy for complexity reduction
%    - The need for a rigorous functional-analytic and geometric foundation
%
% 2. Overview of the four source papers and how they fit together
%    - [FH2012]  On Minimal Subspaces in Tensor Representations
%                (Falcó & Hackbusch, Found. Comput. Math. 12, 2012)
%    - [FHN2019] On the Dirac-Frenkel Variational Principle on Tensor Banach Spaces
%                (Falcó, Hackbusch & Nouy, Found. Comput. Math. 19, 2019)
%    - [FHN2021] Tree-based Tensor Formats
%                (Falcó, Hackbusch & Nouy, SeMA J. 78, 2021)
%    - [FHN2023] Geometry of Tree-based Tensor Formats in Tensor Banach Spaces
%                (Falcó, Hackbusch & Nouy, Ann. Mat. Pura Appl., 2023)
%
% 3. What this monograph adds beyond the articles
%    - Unified notation and language
%    - Expanded proofs and additional examples
%    - New introductory and contextual material
%    - Comparison with parallel literature
%
% 4. Reader's guide: prerequisites, chapter dependencies
%
% 5. Acknowledgements

%------------------------------------------------------------
\chapter*{Notation and Conventions}
\label{chap:notation}
%------------------------------------------------------------
% [~6 pp.]
%
% Systematic reference table for all symbols introduced in
% the macro section above.  To be typeset as a two-column table.
%
% SECTIONS:
%
% N.1  General conventions
%      - Field K = R (default) or C
%      - D = {1,...,d}: dimension index set; |D| = d
%      - Alpha, beta, gamma: generic non-empty subsets of D
%      - alpha^c = D \ alpha: complement
%      - 2^D: power set; cardinality denoted #(alpha)
%      - N_0 = N cup {0}: non-negative integers
%
% N.2  Vector spaces and tensor products
%      - V_j (j in D): component normed spaces
%      - V_alpha = _a otimes_{j in alpha} V_j  (algebraic)
%      - V_D = _a otimes_{j in D} V_j         (algebraic tensor space)
%      - ||.||_j: norm on V_j
%      - ||.||_D: norm on V_D  (crossnorm or reasonable crossnorm)
%      - Completion: \bar{V}_{D,||.||_D} = topological tensor Banach space
%      - _a\otimes: algebraic tensor product (finite sums of elementary tensors)
%      - \otimes_{||.||}:  topological tensor product (completion)
%
% N.3  Subspaces and complemented subspaces
%      - G(X): Grassmannian of X (all split subspaces)
%      - G_r(X): Grassmannian of r-dimensional subspaces  (r finite)
%      - P_{U oplus W}: projection onto U along W
%      - L(X,Y): space of linear maps; L(X,Y) for normed, L_a(X,Y) algebraic
%      - GL(X): invertible bounded linear maps on X
%
% N.4  Minimal subspaces
%      - U^min_alpha(v): unique minimal subspace of v in V_alpha
%      - r_alpha(v) = dim U^min_alpha(v): alpha-rank of v
%      - AD(V_D): set of admissible Tucker ranks
%      - AD(V_D, T_D): set of admissible tree-based rank tuples
%
% N.5  Tucker format
%      - r = (r_1,...,r_d): Tucker rank tuple
%      - M_r(V_D): set of tensors in Tucker format with fixed rank r
%      - Z^(D)(v): tangent subspace at v in M_r(V_D)
%
% N.6  Trees and tree-based formats
%      - T_D: dimension partition tree over D (Definition chap.5)
%      - Root: D;  leaves: L(T_D) = {{j}: j in D}
%      - S(alpha): set of sons (children) of alpha in T_D
%      - level(alpha): level of vertex alpha (root has level 0)
%      - depth(T_D) = max_{alpha} level(alpha)
%      - P_k(T_D): partition of D at level k
%      - r = (r_alpha)_{alpha in T_D}: tree-based rank tuple
%      - FT_r(V_D, T_D): set of tensors with fixed tree-based rank r
%      - FT_{<=r}(V_D, T_D): set with bounded tree-based rank
%      - C^(alpha): transfer tensor at vertex alpha
%      - M_beta(C): matricisation of tensor C in direction beta
%
% N.7  Special trees
%      - T_D^Tucker: Tucker tree (depth 1, S(D) = L(T_D))
%      - T_D^TT: Tensor Train tree (binary, linear)
%      - T_D^HT: Hierarchical Tucker tree (balanced binary)
%
% N.8  Manifold objects
%      - atlas A = {(U_i, phi_i)}
%      - T_v M: tangent space at v
%      - E_{P_D}: model Banach space for Tucker manifold
%      - P_v: metric projection onto T_v M_r(V_D)
%
% N.9  Norms and norm conditions
%      - ||.||_eps: injective norm (smallest reasonable crossnorm)
%      - ||.||_pi: projective norm (largest crossnorm)
%      - (Inj): ||.||_D >= ||.||_eps  (key condition for immersion)
%      - (Cross): ||.||_D is a crossnorm

\tableofcontents

\mainmatter

%%%%%%%%%%%%%%%%%%%%%%%%%%%%%%%%%%%%%%%%%%%%%%%%%%%%%%%%%%%%%
%  PART I
%  ALGEBRAIC AND TOPOLOGICAL FOUNDATIONS
%%%%%%%%%%%%%%%%%%%%%%%%%%%%%%%%%%%%%%%%%%%%%%%%%%%%%%%%%%%%%

\part{Algebraic and Topological Foundations}
\label{part:foundations}

%------------------------------------------------------------
\chapter{Introduction and Motivation}
\label{chap:intro}
%------------------------------------------------------------
% [~25 pp.]

\section{The curse of dimensionality}
\label{sec:intro:curse}
% - Exponential growth of d-variate function spaces
% - Discretisation cost: N^d  vs linear-in-d methods
% - Historical context: Tucker (1966), CP (Hitchcock 1927, Carroll-Chang 1970,
%   Harshman 1970), PARAFAC, HOSVD (De Lathauwer et al. 2000)
% - The role of low-rank approximation

\section{Tensor formats: an informal overview}
\label{sec:intro:overview}
% - Tucker format: subspace representation, core tensor
% - Hierarchical Tucker (HT) format [Hackbusch-Kühn 2009]
% - Tensor Train (TT) format [Oseledets 2011, Vidal 2003]
% - Unifying viewpoint: tree-based Tucker formats
% - Applications: PDEs, quantum chemistry (MCTDH, DMRG),
%   uncertainty quantification, machine learning

\section{Functional-analytic challenges}
\label{sec:intro:challenges}
% - Infinite-dimensional setting: normed vs Banach vs Hilbert
% - Non-linearity of fixed-rank sets
% - Existence of best approximations: not automatic (lack of convexity)
% - Manifold structure: not obvious, key difficulty is complementedness
%   of tangent spaces in Banach (non-Hilbert) spaces

\section{Contributions and structure of the monograph}
\label{sec:intro:structure}
% - Main results summarised informally
% - Chapter-by-chapter guide
% - What requires new work beyond [FH2012,FHN2019,FHN2021,FHN2023]
% - Dependency graph of chapters  (figure)

%------------------------------------------------------------
\chapter{Banach Space Prerequisites}
\label{chap:banach}
%------------------------------------------------------------
% [~30 pp.]
% Self-contained background; expert readers may skip to Ch. 3.

\section{Normed spaces and Banach spaces}
\label{sec:banach:basics}
% - Norms, equivalent norms, completeness
% - Bounded linear maps: L(X,Y), operator norm  ||\cdot||_{Y<-X}
% - Dual space X^*; dual norm; bidual X^**
% - Reflexive Banach spaces; uniform convexity; smooth norms
% - Key examples: L^p(I), H^{1,p}(I), Sobolev spaces

\section{Complemented subspaces and projections}
\label{sec:banach:complemented}
% - Split subspaces: U in G(X)  iff  U closed and has closed complement W
%   with X = U oplus W
% - Projection P_{U oplus W}: definition, continuity, norm estimate
% - Proposition: every finite-dimensional subspace is complemented
% - Proposition: in Hilbert spaces every closed subspace is complemented
% - Counterexamples in non-Hilbert Banach spaces
% - The Grassmannian G(X) and G_r(X): set of split / r-dim subspaces

\section{The Grassmann--Banach manifold}
\label{sec:banach:grassmann}
% - G_r(X) as a Banach manifold modelled on L(U,W) for U in G_r(X)
% - Local charts: Psi_U: G(W, X) -> L(U,W)  via
%     Psi_U(U') = P_{W oplus U} |_{U'} o (P_{U oplus W}|_{U'})^{-1}
% - Transition maps and C^infty structure
% - The lemma on common complements (Char_Projections)
% - Proposition (niceinverse): L_{(U,W)}(X,X) is a nilpotent sub-algebra;
%     exp(L) = id + L
% - G(X) as a Banach manifold NOT modelled on a single space

\section{Banach manifolds: definitions and examples}
\label{sec:banach:manifolds}
% - Atlas, charts, compatibility, C^p and analytic atlases
% - Banach manifold modelled on X; Hilbert manifolds
% - Morphisms (differentiable maps between manifolds)
% - Tangent space T_v M; tangent bundle
% - Immersed and embedded submanifolds
% - Example: GL(X) as open subset of L(X,X)
% - Example: any Banach space as manifold modelled on itself

\section{Multilinear maps and Fréchet differentiability}
\label{sec:banach:multilinear}
% - Multilinear maps between normed spaces; boundedness <-> continuity
% - Proposition: every continuous multilinear map is C^infty-Fréchet
%   differentiable (D^k M = 0 for k > d)
% - Analytic maps over complex Banach spaces:
%   analytic <-> C^1-Fréchet (Proposition analytic)
% - Corollary: continuous multilinear maps between complex Banach spaces
%   are analytic

%------------------------------------------------------------
\chapter{Algebraic Tensor Spaces}
\label{chap:algtensor}
%------------------------------------------------------------
% [~25 pp.]

\section{The algebraic tensor product}
\label{sec:algtensor:def}
% - Universal property and construction (Greub)
% - Elementary tensors v_1 otimes ... otimes v_d
% - Algebraic tensor space V_D = _a otimes_{j in D} V_j
% - Algebraic tensor of linear maps A = otimes A_j
% - Algebraic dual L(V,K); algebraic tensor of dual functionals

\section{Multilinear rank and matricisation}
\label{sec:algtensor:rank}
% - alpha-matricisation M_alpha(v): V_D -> V_alpha otimes V_{alpha^c}
%   as matrix  (finite-dim case: unfolding)
% - rank_j(v) = rank(M_j(v)): Hitchcock rank [1927]
% - Multilinear (Tucker) rank: tuple r = (r_1,...,r_d)
% - Relations between marginal ranks: dim U^min_alpha = dim U^min_{alpha^c}

\section{Partitions and iterated tensor products}
\label{sec:algtensor:partitions}
% - Non-trivial partitions P_D of D
% - Canonical identification: V_D = _a otimes_{alpha in P_D} V_alpha
% - Iterated tensor products; associativity at algebraic level
% - Key example: V_{1,2,3,4} = (V_1 otimes V_2) otimes (V_3 otimes V_4)

\section{The Tucker format}
\label{sec:algtensor:tucker}
% - Tucker subspace representation: v = sum_I C_I prod_j u_j^(i_j)
%   with U_j = span{u_j^(i_j)} in G_{r_j}(V_j)
% - Set T_r: all tensors representable with fixed multi-rank r
% - Equivalent description via subspaces (no explicit basis needed)
% - V_D = union_{r in AD(V_D)} M_r(V_D): disjoint decomposition

%------------------------------------------------------------
\chapter{Minimal Subspaces in Tensor Representations}
\label{chap:minsubspaces}
%------------------------------------------------------------
% [~40 pp.]
% Expanded version of [FH2012]

\section{Definition and existence}
\label{sec:min:def}
% - Definition of U^min_alpha(v) via universality:
%     v in U^min_alpha otimes V_{alpha^c}  and
%     [v in U_alpha otimes V_{alpha^c}  => U^min_alpha subset U_alpha]
% - Theorem: U^min_alpha(v) exists, is unique, has finite dimension
% - Algebraic characterisation:
%   U^min_alpha(v) = span{ (id_alpha otimes phi^(alpha^c))(v) : phi in V^*_{alpha^c} }
% - Symmetry: dim U^min_alpha(v) = dim U^min_{alpha^c}(v)

\section{Lattice properties}
\label{sec:min:lattice}
% - Monotonicity: if v in U_alpha otimes V_{alpha^c} then U^min_alpha(v) subset U_alpha
% - Intersection lemma: (X_1 otimes X_2) cap (Y_1 otimes Y_2) = (X_1 cap Y_1) otimes (X_2 cap Y_2)
% - Inclusion under refinement of partitions (Proposition inclusion_Umin):
%     U^min_alpha(v) subset _a otimes_{beta in P_alpha} U^min_beta(v)
% - New characterisation of minimal subspaces (Prop. (Ualpha in Tensor Uj)):
%   U^min_beta(v) via contractions of U^min_alpha(v)

\section{The alpha-rank and admissible ranks}
\label{sec:min:admissible}
% - alpha-rank: r_alpha(v) = dim U^min_alpha(v)
% - Tucker rank of v: tuple r = (r_j(v))_{j in D} in N_0^d
% - Admissible Tucker ranks: AD(V_D) = {r(v): v in V_D}
% - Characterisation of admissible tuples
% - The set M_r(V_D) for r in AD(V_D): tensors of fixed Tucker rank r
% - V_D = disjoint union over r in AD(V_D) of M_r(V_D)

\section{Best approximation in Banach tensor spaces}
\label{sec:min:bestapprox}
% - Topological tensor product and crossnorms
%   (preliminary; full treatment in Ch. 5)
% - Weak closedness of M_{<=r}(V_D) = FT_{<=r}(V_D, T_D^Tucker)
% - Theorem [FH2012]: in a tensor Banach space with ||.||_D >= ||.||_eps,
%   for any v in V_D there exists best approximation in M_{<=r}(V_D)
% - Extension to intersections of finitely many tensor Banach spaces
% - Examples with Sobolev spaces H^{1,p}(I_1) hat_otimes H^{1,p}(I_2)

\section{Historical notes and comparison}
\label{sec:min:history}
% - Original results in [FH2012]
% - Relation to Schmidt (d=2), SVD, Tucker (1966)
% - Comparison with Falcó-Nouy convex optimization results

%%%%%%%%%%%%%%%%%%%%%%%%%%%%%%%%%%%%%%%%%%%%%%%%%%%%%%%%%%%%%
%  PART II
%  TREE-BASED TENSOR FORMATS
%%%%%%%%%%%%%%%%%%%%%%%%%%%%%%%%%%%%%%%%%%%%%%%%%%%%%%%%%%%%%

\part{Tree-Based Tensor Formats}
\label{part:tree}

%------------------------------------------------------------
\chapter{Dimension Partition Trees}
\label{chap:trees}
%------------------------------------------------------------
% [~20 pp.]

\section{Definition and basic properties}
\label{sec:trees:def}
% - Definition of dimension partition tree T_D (Definition partition_tree):
%   (a) all vertices are non-empty subsets of D
%   (b) D is the root
%   (c) every non-leaf vertex alpha has sons S(alpha) = non-trivial partition of alpha
%   (d) leaves = singletons = L(T_D) = {{j}: j in D}
% - Level of a vertex; depth of T_D
% - Partitions at each level: P_k(T_D) for 0 <= k <= depth(T_D)
%   with P_0 = {D}, P_{depth} = L(T_D) = trivial partition

\section{Special trees}
\label{sec:trees:special}
% - Tucker tree T_D^Tucker: depth 1, S(D) = L(T_D)
%   (Figure: star tree for D={1,2,3,4,5,6})
% - Binary trees; balanced binary trees
% - Hierarchical Tucker tree T_D^HT: balanced binary, depth ~ log_2(d)
% - Tensor Train tree T_D^TT: linear/chain tree, depth = d-1
%   (Figure: chain tree for D={1,...,6})
% - General non-binary trees (Example: D={1,2,3,4,5,6}, depth=2)
%   (Figure from [FHN2021])

\section{Dimension partition tree as a graph}
\label{sec:trees:graph}
% - T_D as rooted tree (graph-theoretic viewpoint)
% - Number of vertices: #T_D
% - Number of internal vertices: #T_D - d
% - Recursive structure: T_D determines T_alpha for each alpha in T_D
% - Isomorphisms of trees; canonical forms

%------------------------------------------------------------
\chapter{Algebraic Tree-Based Tensor Formats}
\label{chap:algTBT}
%------------------------------------------------------------
% [~35 pp.]
% Expanded version of [FHN2021], algebraic part

\section{Tree-based tensor spaces}
\label{sec:algTBT:spaces}
% - Pair (V_D, T_D) as a representation in tree-based format
% - Associated spaces {V_alpha}_{alpha in T_D \ D}
% - For each alpha in T_D \ L(T_D):  V_alpha = _a otimes_{beta in S(alpha)} V_beta
% - Canonical identifications with V_D via iterated tensor products

\section{Minimal subspaces compatible with a tree}
\label{sec:algTBT:mincompat}
% - For v in V_D and alpha in T_D, U^min_alpha(v) is well-defined
% - Compatibility condition: U^min_alpha(v) subset _a otimes_{beta in S(alpha)} U^min_beta(v)
%   (Proposition inclusion_Umin applied at each level of T_D)
% - New Proposition (Prop. (Ualpha in Tensor Uj) for trees):
%   U^min_beta(v) = span{ (id_beta otimes phi^(alpha\beta))(v_alpha) :
%     v_alpha in U^min_alpha(v), phi in ...}

\section{Tree-based rank}
\label{sec:algTBT:tbrank}
% - Tree-based rank: rank_{T_D}(v) = (dim U^min_alpha(v))_{alpha in T_D}
%   in N_0^{#T_D}
% - Admissible tuples: AD(V_D, T_D)
% - Reduction to Tucker rank when T_D = T_D^Tucker
% - Relations between tree-based ranks for different trees

\section{Sets of tensors with fixed and bounded tree-based rank}
\label{sec:algTBT:sets}
% - FT_r(V_D, T_D): fixed rank r in AD(V_D, T_D)
% - FT_{<=r}(V_D, T_D): bounded rank
% - Theorem: FT_{<=r}(V_D, T_D) = union over s<=r of FT_s(V_D, T_D)
% - Relation to Tucker manifolds at each level:
%   FT_r(V_D, T_D) = intersection_{k=1}^{depth(T_D)} M_{r_k}(V_{P_k(T_D)})
%   (Main structural theorem, Theorem 3.12 in [FHN2023])

\section{The transfer tensor representation}
\label{sec:algTBT:transfer}
% - Theorem characterisation_FT [FHN2021]:
%   v in FT_r(V_D,T_D) <=>
%   there exist bases {u^(k)_{i_k}} of U^min_k(v) (k leaf)
%   and transfer tensors C^(alpha) in R^{r_alpha x prod_{beta in S(alpha)} r_beta}
%   with rank conditions rank(M_alpha(C^(alpha))) = r_alpha
%   such that v is expressed recursively via (TBRT1) and (TBRT2)
% - Full proof (expanded from [FHN2021])
% - Special case: Tucker format (one level, C^(D) = core tensor G)
% - Special case: Tensor Train (MPS representation)
% - Special case: Hierarchical Tucker (recursive tree structure)

\section{Matricisation and rank conditions}
\label{sec:algTBT:matricisation}
% - Definition of M_beta: R^{prod r_mu} -> R^{r_beta x prod_{mu!=beta} r_mu}
% - Rank condition: rank(M_beta(C^(alpha))) = r_beta for beta in S(alpha)
%   <=> C^(alpha) in M_r(R^{prod r_mu})  (full-rank condition)
% - Consequences for uniqueness of the representation

%------------------------------------------------------------
\chapter{Topological Tree-Based Tensor Spaces}
\label{chap:topTBT}
%------------------------------------------------------------
% [~25 pp.]
% Expanded from [FHN2021], topological part

\section{Crossnorms and reasonable crossnorms}
\label{sec:topTBT:crossnorms}
% - Crossnorm: ||v_1 otimes ... otimes v_d||_D = prod ||v_j||_j
% - Reasonable crossnorm: both ||.||_D and ||.||_D^* are crossnorms
% - Injective norm ||.||_eps: smallest reasonable crossnorm
% - Projective norm ||.||_pi: largest crossnorm
% - Condition (Inj): ||.||_D >= ||.||_eps  <=> ||.||_D reasonable crossnorm
% - Remark: condition (Inj) is weaker than requiring ||.||_D = ||.||_eps

\section{Topological tree-based tensor spaces}
\label{sec:topTBT:spaces}
% - T_D-continuous tensor product map (Definition)
% - Existence of norms ||.||_alpha for each alpha in T_D \ L(T_D)
% - Theorem ext_Banach [FHN2021]:
%   Under T_D-continuity:  ||.||_alpha-otimes_{beta in S(alpha)} V_{beta,||.||} =
%     ||.||_alpha-otimes_{beta in S(alpha)} V_beta  =  ||.||_alpha-otimes_{j in alpha} V_j
% - Definition: topological tree-based tensor space {V_{alpha,||.||_alpha}}_{alpha in T_D \ {D}}
%   with ambient space V_{D,||.||_D}

\section{Existence of best approximations from tree-based sets}
\label{sec:topTBT:bestapprox}
% - Weak closedness of FT_{<=r}(V_D, T_D) in V_{D,||.||_D}
%   under weakened crossnorm conditions (main improvement over [FH2012])
% - Theorem [FHN2021]: in a topological tree-based tensor space
%   satisfying (Inj) at each level, best approximations exist in FT_{<=r}
% - Proof strategy: induction on depth(T_D), using weak closedness at each level
% - Examples:
%   (a) L^p spaces: ||.||_{0,p} is a reasonable crossnorm for all 1<=p<inf
%   (b) Sobolev spaces H^{1,p}: ||.||_{(0,1),p} is a crossnorm
%   (c) Tree-based Sobolev product spaces (Figure tree for Sobolev example)

%%%%%%%%%%%%%%%%%%%%%%%%%%%%%%%%%%%%%%%%%%%%%%%%%%%%%%%%%%%%%
%  PART III
%  DIFFERENTIAL GEOMETRY OF TENSOR FORMATS
%%%%%%%%%%%%%%%%%%%%%%%%%%%%%%%%%%%%%%%%%%%%%%%%%%%%%%%%%%%%%

\part{Differential Geometry of Tensor Formats}
\label{part:geometry}

%------------------------------------------------------------
\chapter{The Tucker Manifold}
\label{chap:tucker_manifold}
%------------------------------------------------------------
% [~45 pp.]
% Expanded version of [FHN2019], Sections 3-4

\section{The tensor space as a disjoint union}
\label{sec:tucker_manifold:decomp}
% - V_D = disjoint union_{r in AD(V_D)} M_r(V_D)
% - Each M_r(V_D) as a connected component (motivation)
% - Strategy: construct C^inf atlas on each M_r(V_D)

\section{The rank map and the bundle projection}
\label{sec:tucker_manifold:rankmap}
% - Rank map rho: V_D -> union G_{r_alpha}(V_alpha)  given by
%     v |-> (U^min_alpha(v))_{alpha in D}
% - Restricted map rho_r: M_r(V_D) -> prod_{alpha in D} G_{r_alpha}(V_alpha)
%   which is surjective
% - Fibres: rho_r^{-1}(subspace tuple) = M_r(_a otimes U_alpha)
%   is diffeomorphic to M_r(R^{x r_alpha}) (full-rank core tensors)
% - The product manifold prod G_{r_j}(V_j) as base

\section{Local charts for M_r(V_D)}
\label{sec:tucker_manifold:charts}
% - At each v in M_r(V_D):
%     choose W^min_alpha(v) closed complement of U^min_alpha(v) in V_{alpha,||.||_alpha}
%     i.e. V_{alpha,||.||_alpha} = U^min_alpha(v) oplus W^min_alpha(v)
% - Chart domain: U(v) = {w in M_r(V_D): U^min_alpha(w) in G(W^min_alpha(v), V_alpha)}
% - Chart map xi_v: U(v) -> prod L(U^min_alpha(v), W^min_alpha(v)) x ...
%     w |-> (Psi^(alpha)_v(U^min_alpha(w)))_{alpha in D}  x  [core coefficient]
%   where Psi^(alpha)_v = chart map of G_{r_alpha}(V_alpha) at U^min_alpha(v)
% - Explicit formula:
%     Psi^(alpha)_v(U) = P_{W^min_alpha(v) oplus U^min_alpha(v)} |_U
%       circ (P_{U^min_alpha(v) oplus W^min_alpha(v)} |_U)^{-1}  in L(U^min_alpha(v), W^min_alpha(v))
% - Model Banach space:
%     E_{P_D} = prod_{alpha in D} L(U^min_alpha(v), W^min_alpha(v))
%                 x R^{x_{alpha in D} r_alpha}_{*}  (full-rank core tensors)

\section{The C^inf-Banach manifold structure}
\label{sec:tucker_manifold:cinfty}
% - Theorem Tucker_Banach_Manifold [FHN2019]:
%   Under condition (Cross) [tensor product map is continuous]:
%   M_r(V_D) is a C^inf-Banach manifold modelled on E_{P_D}
%   (Theorem restated for partitions via Theorem Tucker_Banach_ManifoldXXX)
% - Full proof with transition map computation
% - Remark: the geometric structure is independent of the choice of ||.||_D
%   (depends only on topology of V_{alpha,||.||_alpha})
% - Example: M_r(_a otimes L^p(I_alpha)) as Banach manifold

\section{The tangent space}
\label{sec:tucker_manifold:tangent}
% - T_v M_r(V_D) = Z^(D)(v) as a closed subspace of V_{D,||.||_D}
% - Characterisation:
%     Z^(D)(v) = sum_{alpha in D} W^min_alpha(v) otimes _a U^min_{D\{alpha}}(v)
%                + _a otimes_{alpha in D} U^min_alpha(v)
% - The tangent space as a sum of d+1 subspaces (geometric interpretation)
% - Proof that Z^(D)(v) is closed in V_{D,||.||_D}  [key difficulty: Banach, not Hilbert]

\section{Immersion in the ambient Banach space}
\label{sec:tucker_manifold:immersion}
% - Lemma (Wmin tensor Umin is split) [FHN2019, Lemma 4.13]:
%     W^min_alpha(v) atens U^min_{D\{alpha}}(v) in G(V_{D,||.||_D})
%   under condition (Inj)
% - Lemma (sum of split subspaces) [FHN2019, Lemma 4.14]
% - Theorem 4.15 [FHN2019]: under condition (Inj),
%     Z^(D)(v) in G(V_{D,||.||_D})  for all v in M_r(V_D)
%   hence M_r(V_D) is an immersed submanifold of V_{D,||.||_D}
% - Corollary: in reflexive V_{D,||.||_D}, best projection onto Z^(D)(v) exists
% - Examples with L^p (Example 4.16) and H^{1,p} (Example 4.17)

%------------------------------------------------------------
\chapter{The Tree-Based Manifold}
\label{chap:TB_manifold}
%------------------------------------------------------------
% [~40 pp.]
% Expanded version of [FHN2023]

\section{Intersection of Tucker manifolds}
\label{sec:TB_manifold:intersection}
% - Key structural result (Theorem 3.12 of [FHN2023] = Thm in manifold paper):
%     FT_r(V_D, T_D) = intersection_{k=1}^{depth(T_D)} M_{r_k}(V_{P_k(T_D)})
%   where r_k = (r_alpha)_{alpha in P_k(T_D)}
% - Each M_{r_k}(V_{P_k}) is a Tucker manifold for the partition P_k
% - The intersection as a submanifold of the ambient space
% - Degenerate cases: when the intersection collapses to a Tucker manifold
%   (Example with depth-2 tree and "trivial" upper level)

\section{Charts and atlas for FT_r(V_D, T_D)}
\label{sec:TB_manifold:charts}
% - At each v in FT_r(V_D, T_D):
%     chart domain compatible with all levels simultaneously
%   U(v) = intersection_{k=1}^{depth} U_k(v)
%   where U_k(v) = chart domain of M_{r_k}(V_{P_k}) at v
% - Chart map: product of chart maps from each Tucker level
%     xi_v = (xi^{(k)}_v)_{k=1}^{depth}
% - Model Banach space:
%     E_{T_D} = prod_{k=1}^{depth} E_{P_k}  (product of Tucker model spaces)
%             modulo compatibility conditions
% - Compatibility with Tucker atlas [FHN2023, Section 4]:
%   the atlas of FT_r restricts to the atlas of M_{r_k} at each level

\section{The C^inf-Banach manifold structure}
\label{sec:TB_manifold:cinfty}
% - Theorem [FHN2023, main theorem]:
%   Under T_D-continuity of tensor product map:
%   FT_r(V_D, T_D) is a C^inf-Banach manifold immersed in V_{D,||.||_D}
% - Proof: induction on depth(T_D) using Tucker manifold results
% - Natural character of the description: comparison with
%   (a) Uschmajew-Vandereycken [2012]: HT format, finite-dim quotient manifold
%   (b) Holtz-Rohwedder-Schneider [2012]: TT format, finite-dim
%   (c) Falcó-Hackbusch-Nouy [2015, arXiv]: different chart system
%   The present description is more natural and compatible with Tucker

\section{Topology and connected components}
\label{sec:TB_manifold:topology}
% - Are the sets FT_r(V_D,T_D) connected? (open question in general)
% - Known results for Tucker manifolds (connected for r != 0)
% - Stratification of V_D by tree-based rank:
%     V_{D,||.||_D} = union_{r in AD(V_D,T_D)} FT_r(V_D,T_D)
% - Boundary and closure: FT_r is NOT closed; FT_{<=r} IS weakly closed (Ch. 7)
% - Open problems: global topology, fundamental groups

%%%%%%%%%%%%%%%%%%%%%%%%%%%%%%%%%%%%%%%%%%%%%%%%%%%%%%%%%%%%%
%  PART IV
%  VARIATIONAL PRINCIPLES AND APPLICATIONS
%%%%%%%%%%%%%%%%%%%%%%%%%%%%%%%%%%%%%%%%%%%%%%%%%%%%%%%%%%%%%

\part{Variational Principles and Applications}
\label{part:variational}

%------------------------------------------------------------
\chapter{The Dirac--Frenkel Variational Principle}
\label{chap:DF}
%------------------------------------------------------------
% [~35 pp.]
% Expanded from [FHN2019] Sect. 5, and [FHN2023] Sect. 5

\section{The abstract Cauchy problem on tensor Banach spaces}
\label{sec:DF:cauchy}
% - Abstract ODE: du/dt = F(t,u(t)),  u(0) = u_0  in V_{D,||.||_D}
% - Assumptions on F: Lipschitz, existence and uniqueness of solutions
% - Goal: approximate solution by a curve t |-> v_r(t) in M_r(V_D)  [or FT_r]
%   (reduced order model / model reduction)
% - Physical motivation: time-dependent Schrödinger equation,
%   MCTDH (Multi-Configuration Time-Dependent Hartree),
%   stochastic parametric PDEs

\section{The Dirac--Frenkel principle: Hilbert space case}
\label{sec:DF:hilbert}
% - Classical formulation in Hilbert space H:
%     <v_dot_r(t) - F(t,v_r(t)), v_dot> = 0  for all v_dot in T_{v_r(t)} M_r
%   <=> v_dot_r(t) = P_{v_r(t)} F(t, v_r(t))
%   where P_v = metric projection onto T_v M_r (Hilbert projection)
% - Well-posedness in finite dimension (Koch-Lubich 2007)
% - Limitations: Hilbert structure needed for metric projection

\section{Metric projection and convexity in Banach spaces}
\label{sec:DF:projection}
% - Strictly convex Banach spaces: unique metric projections onto convex sets
% - Uniformly convex spaces (Clarkson): reflexive + strictly convex
% - Smooth Banach spaces: Fréchet differentiability of the norm
% - Metric projection onto Z^(D)(v):
%     v_dot_r(t) = argmin_{v_dot in Z^(D)(v_r(t))} ||v_dot - F(t,v_r(t))||_D
%   existence guaranteed when V_{D,||.||_D} is reflexive (Corollary 4.18 [FHN2019])
%   uniqueness when V_{D,||.||_D} is strictly convex

\section{Extension to Tucker format: Dirac--Frenkel in Banach spaces}
\label{sec:DF:tucker}
% - Theorem 5.2 [FHN2019]:  (existence of reduced ODE on Tucker manifold)
%   Under conditions (Cross) and (Inj):
%   for any smooth vector field F on V_{D,||.||_D},
%   there exists a solution of the projected ODE on M_r(V_D)
%   (at least locally in time)
% - Theorem 5.4 [FHN2019]:  (characterisation via variational condition)
%   v_dot_r(t) in Z^(D)(v_r(t)) and
%   ||v_dot_r(t) - F(t,v_r(t))||_D = min_{v_dot in Z^(D)(v_r(t))} ||v_dot - F||_D
% - Full proof (expanded from [FHN2019])
% - Remark: in Hilbert case recovers classical Dirac-Frenkel equation

\section{Extension to tree-based format}
\label{sec:DF:treebased}
% - Theorem [FHN2023, Section 5]: Theorems 5.2 and 5.4 of [FHN2019]
%   extend verbatim to FT_r(V_D, T_D) when:
%   (a) Condition (Inj) holds
%   (b) The tensor product map is T_D-continuous
% - Proof: use the immersion result (Ch. 9) and note that
%   T_{v} FT_r(V_D,T_D) is a closed split subspace of V_{D,||.||_D}
% - Consequences for DMRG and ALS algorithms (informal discussion)

%------------------------------------------------------------
\chapter{Open Problems and Perspectives}
\label{chap:open}
%------------------------------------------------------------
% [~20 pp.]

\section{Algorithmic implications}
\label{sec:open:algorithms}
% - ALS (Alternating Least Squares) and its geometric interpretation
%   on M_r(V_D) or FT_r(V_D,T_D)
% - DMRG (Density Matrix Renormalisation Group) as optimisation on TT manifold
% - Riemannian gradient methods on tensor manifolds (Uschmajew-Vandereycken 2020)
% - Convergence analysis relying on manifold geometry (curvature, injectivity radius)

\section{Global topology of tensor manifolds}
\label{sec:open:topology}
% - Connectedness of FT_r(V_D,T_D): proven for Tucker (r != 0), open for general trees
% - Higher homotopy groups of Grassmannians and Tucker manifolds
% - Stratification by rank: FT_r as stratum, FT_{<=r} = closure (Zariski-like)
% - Degree and characteristic classes: analogues of Schubert calculus?

\section{Tensor formats over non-classical spaces}
\label{sec:open:spaces}
% - Extension to Fréchet spaces, locally convex spaces
% - Manifolds of tensors in quasi-Banach spaces (L^p for 0 < p < 1)
% - Infinite-dimensional index sets: nuclear spaces, projective limits

\section{Learning on tensor manifolds}
\label{sec:open:learning}
% - Tensor networks in machine learning: MPS, MERA, PEPS
% - Gradient descent on FT_r as optimisation on a Banach manifold
% - Generalisation bounds via geometry (tube formula, Weyl)
% - Open: role of tree choice in approximation quality (optimal tree selection)

\section{Connections with quantum information}
\label{sec:open:quantum}
% - Matrix Product States (MPS) = TT format in C^N (complex Hilbert space)
% - Entanglement entropy and Schmidt ranks = tree-based ranks
% - Area laws and tree tensor network states
% - Relation to C^*-algebras and operator systems

%%%%%%%%%%%%%%%%%%%%%%%%%%%%%%%%%%%%%%%%%%%%%%%%%%%%%%%%%%%%%
%  APPENDICES
%%%%%%%%%%%%%%%%%%%%%%%%%%%%%%%%%%%%%%%%%%%%%%%%%%%%%%%%%%%%%

\appendix

%------------------------------------------------------------
\chapter{Banach Manifolds: A Summary}
\label{app:manifolds}
%------------------------------------------------------------
% [~8 pp.]
% Quick reference for the main definitions and theorems used.
% No proofs; references given to Lang [Fundamentals of Differential Geometry]
% and Abraham-Marsden-Ratiu [Manifolds, Tensor Analysis, Applications].
%
% A.1  Topological manifolds
% A.2  Banach manifolds; atlas; charts; C^p and analytic structures
% A.3  Morphisms; diffeomorphisms; local diffeomorphism theorem
% A.4  Tangent space; tangent bundle; differential
% A.5  Submanifolds: immersed vs embedded
% A.6  Vector fields; flows; ODE on manifolds
% A.7  Riemannian and Finsler metrics on Banach manifolds (brief)

%------------------------------------------------------------
\chapter{Functional Analysis: Selected Results}
\label{app:functional}
%------------------------------------------------------------
% [~5 pp.]
% A.1  Hahn-Banach theorem and corollaries
% A.2  Open mapping theorem; closed graph theorem
% A.3  Uniform boundedness principle
% A.4  Reflexive spaces; weak convergence; Eberlein-Šmulian theorem
% A.5  Proximinal sets; best approximation in reflexive spaces
% A.6  Strictly and uniformly convex spaces: key properties

%------------------------------------------------------------
\chapter{Index of Notation}
\label{app:notation_index}
%------------------------------------------------------------
% [~3 pp.]
% Alphabetical list of all symbols with page reference.
% Cross-references to the Notation chapter (front matter).

\backmatter

\bibliographystyle{alpha}
\bibliography{monograph}

%------------------------------------------------------------
% BIBLIOGRAPHY NOTES (for the author -- not typeset)
%
% Core references to include:
%
% [FH2012]      Falcó-Hackbusch, Found. Comput. Math. 12 (2012) 765-803
% [FHN2019]     Falcó-Hackbusch-Nouy, Found. Comput. Math. 19 (2019) 159-204
% [FHN2021]     Falcó-Hackbusch-Nouy, SeMA J. 78 (2021) 159-173
% [FHN2023]     Falcó-Hackbusch-Nouy, Ann. Mat. Pura Appl. (2023)
%
% Tensor formats:
% [Hackbusch]   Hackbusch, Tensor Spaces and Numerical Tensor Calculus,
%               2nd ed., Springer 2019
% [HaKuehn2009] Hackbusch-Kühn, J. Fourier Anal. Appl. 15 (2009) 706-722
% [Osedelets1]  Oseledets, SIAM J. Sci. Comput. 33 (2011) 2295-2317
% [Vidal]       Vidal, Phys. Rev. Lett. 91 (2003) 147902
%
% Geometry of tensor formats:
% [KoLu1]       Koch-Lubich, SIAM J. Matrix Anal. Appl. 29 (2007) 434-454
% [Uschmajew2012] Uschmajew-Vandereycken, Lin. Alg. Appl. 437 (2012)
% [HRS]         Holtz-Rohwedder-Schneider, Found. Comput. Math. 12 (2012) 
% [AMS]         Absil-Mahony-Sepulchre, Optimization Algorithms on Matrix
%               Manifolds, Princeton UP 2008
%
% Functional analysis / Banach geometry:
% [Lang]        Lang, Fundamentals of Differential Geometry, Springer 1999
% [Greub]       Greub, Linear Algebra, 4th ed., Springer 1981
% [Douady66]    Douady, Bull. Soc. Math. France 94 (1966) 323-374
% [FHHM]        Falcó-Hackbusch-Nouy... (Grassmannian results)
% [DRW]         Douady-Robba-Wallach? (check [FHN2019] reference list)
%
% Applications:
% [Bachmayr2016] Bachmayr et al., ESAIM Math. Model. Numer. Anal. (2016)
% [nouy:2017]   Nouy, in: Model Reduction and Approximation, SIAM 2017
% [Lubish]      Lubich, From Quantum to Classical Molecular Dynamics, EMS 2008
% [KOL09]       Kolda-Bader, SIAM Rev. 51 (2009) 455-500
%------------------------------------------------------------

\printindex

\end{document}
%============================================================
%  END OF SKELETON
%============================================================
